\section{Introduction}

The Cocoon project targets a wireless access paradigm in which the network adapts to the fluidity of the city rather than assuming a static infrastructure and a stationary demand profile. The motivating hypothesis from the project proposal is that a practical \gls{cf} network becomes more feasible when low-height infrastructure can be co-created, repositioned, and activated only when the measured channel conditions justify it. In that view, the citizens are not only traffic sources but also anchors for sensing, mounting, and reconfiguring the surrounding network.

This article narrows that broader vision to a simulation-based study of low-height infrastructure adaptation. The implemented framework covers wall-mounted candidate \glspl{ap}, outdoor user mobility, path-based \gls{csi} extraction, and repeated placement decisions over relocation windows in an urban scene. This scope enables a reproducible comparison of geometry, radio-performance, and relocation-cost metrics for alternative infrastructure strategies.

\PaperFigure[t][0.86\linewidth]{\ConceptFigurePath}{Conceptual illustration of the Cocoon low-height access-network vision, where citizens interact with the network and move \glspl{ap} depending on the current channel conditions and its users. By dynamically adapting to the current network state and requirement, the ecological best solution can be achieved when all network resources can cooperate in terms of energy efficiency and minimum materials needed. The lines between the people represent the real-time channel gain measurements. Each person thereby represents a potential \gls{ap} candidate position. Based on these measurements, the most optimal \gls{ap} location can be determined, which is in this case the person close to the tree.}{fig:cocoon_concept}

The main contributions of the present study are:
\begin{itemize}
    \item a concise technical framing of Cocoon as a low-height, movable-\gls{ap}, citizen-co-created access-network concept,
    \item a reproducible description of the Rabot outdoor simulation pipeline implemented in \texttt{COCOON-Sionna}, and
    \item a common analysis workflow for geometry, \gls{sinr}, \gls{esr}, and relocation-cost results derived from exported simulation artifacts.
\end{itemize}
